\chapter{Iterators, Generators, and \texttt{yield from} (Python 3.3)}

This chapter is the canonical home for the \textbf{iteration model} that pervades Python ---
\texttt{for} loops, comprehensions, unpacking, \texttt{in}, and ultimately \texttt{async for} all
rest on it. We build it from the bottom: the \textbf{iterator protocol}
(\texttt{\_\_iter\_\_}/\texttt{\_\_next\_\_}), then \textbf{generators} as the ergonomic way to
write iterators by suspending a frame, then \textbf{\texttt{yield from}} (PEP 380), the delegation
operator that turned generators into composable coroutines and laid the last stone before native
\texttt{async}/\texttt{await} (Chapter 11). We close with \textbf{implicit namespace packages}
(PEP 420). Chapter 5 introduced generator \texttt{send}/\texttt{throw}/\texttt{close}; here we
complete the model.

\section*{Section Index}
\begin{itemize}
  \item \textbf{9.1} The iterator protocol
  \item \textbf{9.2} Generators: iterators written as suspended frames
  \item \textbf{9.3} \texttt{yield from}: delegation and the coroutine bridge (PEP 380)
  \item \textbf{9.4} Implicit namespace packages (PEP 420)
  \item \textbf{9.5} Performance and anti-patterns
  \item \textbf{9.6} Summary and cross-references
\end{itemize}

\section{The iterator protocol}

\textbf{Why this exists.} Python decouples \emph{consuming} a sequence from \emph{being} a
sequence, so one \texttt{for} loop drives a list, a file, a cursor, an infinite counter, or a
network stream. An \textbf{iterable} implements \texttt{\_\_iter\_\_()} (returns a fresh
iterator); an \textbf{iterator} implements \texttt{\_\_next\_\_()} (next item or
\texttt{StopIteration}) and \texttt{\_\_iter\_\_()} returning \texttt{self}.

\begin{lstlisting}[language=Python, caption={A hand-written iterator; for-loops speak exactly this protocol.}]
class Countdown:
    def __init__(self, n):
        self.n = n
    def __iter__(self):
        return self
    def __next__(self):
        if self.n <= 0:
            raise StopIteration
        self.n -= 1
        return self.n + 1

print("iterate:", list(Countdown(3)))
\end{lstlisting}

Verified output (CPython 3.13.5):

\begin{lstlisting}[caption={Output}]
iterate: [3, 2, 1]
\end{lstlisting}

\begin{lstlisting}[language=Python, caption={The bytecode of any for-loop.}]
import dis
def loop(xs):
    for x in xs:
        pass
dis.dis(loop)
\end{lstlisting}

Verified output (CPython 3.13.5):

\begin{lstlisting}[caption={Output}]
 14           LOAD_FAST                0 (xs)
              GET_ITER
      L1:     FOR_ITER                 3 (to L2)
              STORE_FAST               1 (x)
              JUMP_BACKWARD            5 (to L1)
      L2:     END_FOR
              POP_TOP
              RETURN_CONST             0 (None)
\end{lstlisting}

\texttt{StopIteration} is \emph{control flow}, not an error---\texttt{FOR\_ITER} catches it to end
the loop. (A stray \texttt{StopIteration} inside a generator body is a bug; PEP 479 converts it to
\texttt{RuntimeError}.)

\textbf{The senior-engineer contrast.} Java separates \texttt{Iterable} from \texttt{Iterator}
with a look-ahead \texttt{hasNext()}. Python has no \texttt{hasNext()}; iteration is
\textbf{exhaustion-by-exception}. C++ iterators are \emph{positions} compared against
\texttt{end()}. Python's is the most minimal---and \textbf{iterators are single-pass and
stateful}.

\section{Generators: iterators written as suspended frames}

\textbf{Why this exists.} Writing the iterator class by hand is boilerplate. A \textbf{generator
function}---any \texttt{def} containing \texttt{yield}---lets you write the logic and have CPython
synthesize the iterator. Calling it returns a generator object; each \texttt{next()} runs to the
next \texttt{yield} and suspends.

\begin{lstlisting}[language=Python, caption={A generator is a lazy iterator; its frame lives on the heap between yields.}]
def gen_squares(n):
    for i in range(n):
        yield i * i

g = gen_squares(4)
print("type:", type(g).__name__, "| has a live frame:", g.gi_frame is not None)
print("values:", list(g))
\end{lstlisting}

Verified output (CPython 3.13.5):

\begin{lstlisting}[caption={Output}]
type: generator | has a live frame: True
values: [0, 1, 4, 9]
\end{lstlisting}

\textbf{What the interpreter actually does.} A generator wraps a \texttt{PyGenObject} holding a
\texttt{gi\_frame}. At each \texttt{yield}, CPython saves the resume offset (\texttt{f\_lasti}) and
stack depth, marks the generator suspended, and returns control---the frame is not freed, it lives
on the heap. \texttt{next()}/\texttt{send()} re-attach and resume. This is what makes generators
\textbf{lazy and $O(1)$ in memory}: \texttt{gen\_squares(10**9)} allocates one frame, not a
billion-element list. The two-way protocol (\texttt{send}/\texttt{throw}/\texttt{close}) was
covered in Chapter 5.

\section{\texttt{yield from}: delegation and the coroutine bridge (PEP 380)}

\textbf{Why this exists.} Before 3.3, a generator delegating to a sub-generator \emph{and}
correctly forwarding \texttt{send}/\texttt{throw}/\texttt{close} and capturing its return value had
to hand-write a fragile loop. \textbf{PEP 380's \texttt{yield from}} makes the delegator a
transparent channel: yielded values pass to the caller; \texttt{send} values pass to the
sub-generator; thrown exceptions raise inside it; \texttt{close} propagates; and the
sub-generator's \textbf{return value becomes the value of the \texttt{yield from} expression} (via
\texttt{StopIteration.value}).

\begin{lstlisting}[language=Python, caption={yield from delegates send/throw and captures the sub-generator's return value.}]
def subgen():
    received = yield "a"
    yield f"got {received}"
    return "subgen-result"

def delegator():
    result = yield from subgen()       # transparent channel; captures the return
    print("delegator captured return:", result)
    yield "after"

d = delegator()
print("next:", next(d))                 # -> "a" (from subgen)
print("send:", d.send("X"))             # "X" routed into subgen -> "got X"
print("next:", next(d))                 # subgen returns; result captured; -> "after"
\end{lstlisting}

Verified output (CPython 3.13.5):

\begin{lstlisting}[caption={Output}]
next: a
send: got X
delegator captured return: subgen-result
next: after
\end{lstlisting}

\begin{lstlisting}[language=Python, caption={yield from composes generators cleanly (recursive flatten).}]
def flatten(nested):
    for item in nested:
        if isinstance(item, list):
            yield from flatten(item)
        else:
            yield item

print("flatten:", list(flatten([1, [2, [3, 4], 5], 6])))
\end{lstlisting}

Verified output (CPython 3.13.5):

\begin{lstlisting}[caption={Output}]
flatten: [1, 2, 3, 4, 5, 6]
\end{lstlisting}

\textbf{The coroutine bridge.} \texttt{yield from} let a generator delegate to another that
performs I/O and yields control upward---the exact pattern early \texttt{asyncio} used
(\texttt{@asyncio.coroutine} + \texttt{yield from}) before 3.5 gave it dedicated syntax. When you
read \texttt{await expr} in Chapter 11, picture \texttt{yield from} with a typed, awaitable-only
channel: the suspension mechanism is the same heap frame.

\section{Implicit namespace packages (PEP 420)}

\textbf{Why this exists.} Through 3.2, a directory was a package only with \texttt{\_\_init\_\_.py},
which blocked distributing one logical package across \textbf{multiple installs/directories}.
\textbf{PEP 420} lets a package span directories with \emph{no} \texttt{\_\_init\_\_.py}: a
\textbf{namespace package}.

\textbf{The finder algorithm.} During \texttt{import foo}, \texttt{PathFinder} scans
\texttt{sys.path}. A \texttt{foo/\_\_init\_\_.py} makes a regular package and stops the search. With
no \texttt{\_\_init\_\_.py} but one or more directories named \texttt{foo}, it \textbf{accumulates
all} of them into a namespace package whose \texttt{\_\_path\_\_} lists them and whose
\texttt{\_\_file\_\_} is \texttt{None}.

\begin{lstlisting}[language=Python, caption={One namespace package 'company' spanning two directories (no company/\_\_init\_\_.py in either).}]
import sys
sys.path.insert(0, "path1")
sys.path.insert(0, "path2")

import company             # no __init__.py anywhere -> namespace package
import company.core        # contributed by path1
import company.extra       # contributed by path2

print("company.__file__:", company.__file__)
print("company.__path__:", list(company.__path__))
print("core:", company.core.VALUE, "| extra:", company.extra.VALUE)
\end{lstlisting}

Verified output (CPython 3.13.5):

\begin{lstlisting}[caption={Output}]
company.__file__: None
company.__path__: ['/private/tmp/.../path2/company', '/private/tmp/.../path1/company']
core: from path1 | extra: from path2
\end{lstlisting}

Two directories contribute submodules to one importable \texttt{company} namespace---how
ecosystems ship plugins (\texttt{zope.*}, \texttt{google.cloud.*}) as separately-installable
distributions. \textbf{Trade-off:} namespace resolution is slower (the finder scans all of
\texttt{sys.path}), and accidentally omitting an \texttt{\_\_init\_\_.py} from a regular package
silently turns it into a namespace package. The complete import machinery is Vol X.

\section{Performance and anti-patterns}

\begin{itemize}
  \item \textbf{Generators trade memory for laziness}---$O(1)$ vs $O(N)$; the right default for
    pipelines/streams. The cost is a per-item Python-level resume; for small materialized data a
    list comprehension can be faster and re-iterable.
  \item \textbf{Iterators are single-pass}---exhausting one leaves it empty; a second pass yields
    nothing. Re-create or materialize. Bites with \texttt{zip}/\texttt{map}/dict views iterated
    twice.
  \item \textbf{Don't \texttt{list()} a generator just to index} if you iterate once; \emph{do}
    materialize when you need \texttt{len}, random access, or reuse.
  \item \textbf{\texttt{return} inside a generator} sets \texttt{StopIteration.value} (for
    \texttt{yield from}); a bare \texttt{StopIteration} raised in generator code becomes
    \texttt{RuntimeError} (PEP 479)---just \texttt{return}.
  \item \textbf{Reach for \texttt{itertools}} before hand-rolling generator plumbing (Vol XII).
  \item \textbf{Namespace-package footgun:} a missing \texttt{\_\_init\_\_.py} silently creates one.
\end{itemize}

\section{Summary and cross-references}

\begin{itemize}
  \item The \textbf{iterator protocol} is \texttt{\_\_iter\_\_} + \texttt{\_\_next\_\_}
    (\texttt{StopIteration} to stop); \texttt{for} compiles to \texttt{GET\_ITER}/\texttt{FOR\_ITER};
    iterators are \textbf{single-pass}.
  \item \textbf{Generators} synthesize iterators from \texttt{yield} by \textbf{suspending a heap
    frame}, giving $O(1)$-memory laziness and a two-way protocol.
  \item \textbf{\texttt{yield from}} (PEP 380) delegates transparently and \textbf{captures the
    return value}---the direct ancestor of \texttt{await}.
  \item \textbf{Implicit namespace packages} (PEP 420) span directories without
    \texttt{\_\_init\_\_.py} (\texttt{\_\_file\_\_ is None}, multi-entry \texttt{\_\_path\_\_}).
\end{itemize}

\textbf{Cross-references.} Generator \texttt{send}/\texttt{throw}/\texttt{close} origins
$\rightarrow$ Chapter 5. Comprehensions vs. generator expressions $\rightarrow$ Chapter 3. PEP 393
flexible strings $\rightarrow$ Chapter 7. The buffer protocol and \texttt{memoryview} $\rightarrow$
Vol IX. \texttt{asyncio} and native \texttt{async}/\texttt{await} $\rightarrow$ Chapter 11. The
complete import machinery $\rightarrow$ Vol X. \texttt{itertools} $\rightarrow$ Vol XII.
